%==============================================================================
% Section 8: Discussion
%==============================================================================
\section{Discussion}
\label{sec:discussion}

\subsection{What $\chi$-universality organises}
\label{sec:discussion:organises}

The results of this paper show that, in the Lorentzian EC+NY reduced
sector, topology dependence is organised by the same scalar $\chi$ across
three levels: Hamiltonian admissibility, P-channel diagnostics, and the
reduced vacuum orbit.

At the admissibility level, the $\EH/\AX/\VT/\MX$ mode pattern is common
to the four topologies (Table~\ref{tab:four_topology_admissibility}), and
the specific values of the topology do not alter the pattern itself.  At
the P-channel level, $\Pform$ is exactly zero independently of topology
(Section~\ref{sec:pchannel:Pform}) and $\QNY$ shows the mode separation
(Section~\ref{sec:pchannel:QNY}); neither has $\chi$-dependence.  Only the
MX diagnostic $\Pint$ distinguishes the topologies, through $\Ctop=-9\chi$
(Section~\ref{sec:pchannel:Pint}).  At the reduced-orbit level, the
auxiliary-shell equation reduces to $\dot q^{2}+\chi=0$
(Section~\ref{sec:orbit:derivation}), and the qualitative classification of
orbits is determined by the sign of $\chi$
(Section~\ref{sec:orbit:sign}).

Taken together, the differences among the four topologies appear in the
reduced sector as differences in the value of the scalar $\chi$, while the
remaining structure (the admissibility pattern, the behaviour of $\Pform$,
and the mode separation) is common across topologies.  We call this
phenomenon $\chi$-controlled topology universality.

\subsection{Relation to the paper~IV response dictionary}
\label{sec:discussion:paper04echo}

The comparison with paper~IV~\cite{paper4} suggests that the four Thurston
geometries used throughout DPPU are not merely a convenient list of
examples.  In paper~IV, these four geometries exhausted four response
types under Euclidean KK Higgsing: none, uniaxial, triaxial, and biaxial.
Here, in a Lorentzian spin-0 / isotropic EC+NY reduction, the same four
geometries appear as structurally distinguished loci of the five-parameter
structure-constant family: the origin ($\Tthree$), a single class-A axis
($\Nilthree$), the isotropic class-A diagonal ($\Sthree$), and a balanced
class-B pair ($\Solthree$).  The $\chi$-universality result therefore
provides a new reading of the same four-geometry response dictionary in a
different reduced sector, rather than a direct identification with the
paper~IV mechanism.

These two results suggest that the same set of Thurston geometries gives a
coherent classificatory structure across different reduced sectors of DPPU
theory, but the elucidation of a common underlying basis is left to
future work.

\subsection{Limitations and future directions}
\label{sec:discussion:limitations}

This paper is restricted to the $\alpha=0$ EC+NY isotropic reduced sector,
and the following extensions are explicitly left to future work.

For $\Nilthree$ and $\Solthree$, this paper uses only the reduced
description in terms of the isotropic scale $q$ and does not derive the
anisotropic EOM with independent scale factors.  For $\Solthree$, the
construction of compact quotients, the classification of spin structures,
and the determination of spectral data are outside the scope of the
present paper.  The absence of real initial data on the $\Sthree$ vacuum
reduced atlas may be resolved by the introduction of matter coupling or an
effective source, but that analysis is outside the scope of the present
paper.  As for the Weyl-source / $\alpha\neq0$ branches, this paper is
restricted to the $\alpha=0$ baseline and does not include a re-evaluation
of the other branches.

For the exact cancellation of $\Pform$ by block orthogonality and the
$\chi$-specialisation of $\PintMX$, we present the results as confirmed by
symbolic computation, and do not proceed to abstract formulation in the
form of formal lemmas and propositions.

\subsection{Closing remarks}
\label{sec:discussion:closing}

In summary, this paper has shown that in the Lorentzian EC+NY reduced
sector the four topologies $\Sthree$, $\Tthree$, $\Nilthree$, $\Solthree$
share a common admissibility pattern, and that the topology dependence is
collapsed onto a single scalar $\chi$ through two independent routes: the
diagnostic identity $\Ctop=-9\chi$ and the reduced vacuum orbit
$\dot q^{2}+\chi=0$.  This $\chi$-controlled topology universality is a
classification result in the reduced sector; generalisation to
anisotropic, matter-coupled, and global extensions is a future task
beyond the scope of the present paper.
